\section{Case Studies}
\label{sec:cases}

\S\ref{sec:framework} developed the framework abstractly. \S\ref{sec:results} reports what the four-language experiment observed. This section connects the two. For each of the four languages tested, we describe the morphological system at the level of organisation relevant to the framework, walk through how the grammar engine decomposes representative word forms, and report the numerical inputs ($|B(L)|$, $H(L)$, $\rho(L)$) that the engine produces from a 20{,}000-sentence sample of the mixed Wikipedia and FineWeb corpora. The four subsections follow the project's canonical order (Mandarin, English, Turkish, Arabic), from the most isolating system to the most templatic, the order in which the framework predicted the rebate would grow, a prediction \S\ref{sec:results} puts to the test and, for three of the four languages, overturns.

\subsection{Mandarin: isolating}
\label{sec:cases-zh}

Mandarin sits at the isolating end of the typological spectrum. Grammar is carried by word order and a small inventory of function words, and inflectional morphology in the European sense is essentially absent: verbs do not conjugate for person or tense, nouns do not decline for case, and most words consist of a single morpheme or a stable two-character compound. Grammatical relations that European languages bundle into suffixes are instead handled by free particles: aspect is marked by post-verbal \emph{le} (\zh{了}, perfective), \emph{zhe} (\zh{着}, durative), and \emph{guo} (\zh{过}, experiential); argument structure is rearranged by the \emph{ba} construction (\zh{把}); modality is carried by pre-verbal auxiliaries such as \emph{hui} (\zh{会}, ``will / be able to'') and \emph{neng} (\zh{能}, ``can'').

But Mandarin carries its own structural signal, hidden at a layer below the word. Every Han character is built around a \emph{Kangxi radical}, the semantic component used to organise the traditional dictionary. The radical encodes a coarse semantic class: \zh{子} marks CHILD / KINSHIP terms, \zh{木} marks TREE / WOOD entities, \zh{水} (and its combining form \zh{氵}) marks WATER / LIQUID, \zh{心} (and \zh{忄}) marks HEART / EMOTION, \zh{言} (and \zh{讠}) marks SPEECH / LANGUAGE, and so on across the 214-radical Kangxi system. This radical-semantic axis is the structural analog of Arabic's \emph{wazn}-semantic-class axis: in both languages, a non-concatenative layer of semantic information sits behind the surface form, and a tokeniser that knows to look at it gains access to a signal that flat byte-level segmentation discards.

A representative decomposition: the word \zh{学习} (\emph{xu\'exi}, ``to study'') decomposes into two characters, each carrying its own radical. \zh{学} (\emph{xu\'e}, ``learn'') is built on radical \zh{子} (CHILD), placing it in the HUMAN\_RELATION semantic class; \zh{习} (\emph{x\'i}, ``practise'') is built on radical \zh{羽} (FEATHERS), placing it in the NATURAL\_OBJECT class. The compound is grammatically a verb, but its grammatical identity sits in syntactic position rather than in any inflectional marker on the form itself; aspectual nuance, when needed, is supplied by a following particle (\zh{了}, \zh{着}, \zh{过}).

On the 20{,}000-sentence Mandarin sample the engine serialises $|B(L)| = 891$ distinct bundle strings, populated by part-of-speech, particle-class, and radical-class combinations rather than by inflection. Two of the engine's numbers for Mandarin must be read with care. The raw $H(L) = 6.28$ bits is anomalously high for an isolating language that marks almost no grammar on the word; it reflects the breadth of the character-and-radical serialisation, not a genuine inflectional load, and tightening the Mandarin bundle definition so that $H$ measures grammar rather than serialisation breadth is noted as engine work in \S\ref{sec:conclusion}. What is linguistically meaningful is the structural recoverability, $\rho(\text{Mandarin}) \approx 0.025$: almost none of what the engine serialises is recoverable from the bare character string, because Mandarin's grammar lives in word order and free particles, not on the word surface. The product $\rho(L)\cdot H(L) \approx 0.16$ is accordingly the lowest of the four, the correct conclusion (Mandarin exposes the least surface grammar) reached in spite of the inflated $H$. The \texttt{radical\_class} stream supplies a separate, lexicon-recoverable semantic axis: 214 Kangxi radicals in principle, with roughly 30 covering most of the corpus mass.

\subsection{English: analytic}
\label{sec:cases-en}

English is the next-most analytic of the four languages, and the typological baseline against which the framework's predictions will feel least surprising. English grammar is carried largely by word order and by function words (articles, prepositions, auxiliaries). Inflectional morphology is limited to a small inventory of suffixes (plural \emph{-s}, possessive \emph{'s}, third-person-singular present \emph{-s}, past \emph{-ed}, progressive \emph{-ing}, comparative \emph{-er}, superlative \emph{-est}), plus a closed set of suppletive and irregular forms (\emph{go/went}, \emph{be/was}, \emph{child/children}) resolved through a lookup list. Derivational morphology is richer, and the engine now treats it as a productive layer rather than a lexicalised annotation: affixes are peeled off the surface form, one at a time, until a known root is reached.

A representative decomposition. The word \emph{researchers} resolves to root \emph{search}, a derivational chain \texttt{[er $\to$ AGENT\_NOUN, re $\to$ REPETITION]}, and feature bundle \texttt{pos=NOUN | num=PL}. Each derivational affix is annotated by the semantic role it contributes, rather than collapsed into a lexicalised stem. The word \emph{went} resolves to the suppletive root \emph{go} with bundle \texttt{pos=VERB | tense=PAST}; the surface string bears no phonological relationship to its root, and the lookup list is the only route by which the underlying form is recovered. The word \emph{runs} is ambiguous on the surface alone between a plural noun and a third-person-singular verb; the engine resolves this through part-of-speech inference and records the residual ambiguity with an \texttt{ambig\_3sg=YES} tag.

The observed feature-bundle inventory remains small. The 20{,}000-sentence English sample yields $|B(L)| = 21$ distinct bundles. The corpus-weighted grammatical entropy is $H(L) = 1.30$ bits per word, with upper bound $H_{\max}(L) = \log_2 21 \approx 4.39$ bits. The gap between $H(L)$ and $H_{\max}(L)$ reflects the strong peakedness of English's bundle distribution: most English word forms are uninflected singular nouns or bare verbs, and the entropy of their features is correspondingly low.

Surface ambiguity in English is localised and well-bounded. The \emph{-s} homograph between noun plural and verb 3SG, identified through corpus statistics, is the main contributor to $H(b \mid \text{signature})$. The signature-conditional entropy comes to $0.23$ bits; $\rho(\text{English}) = 1 - 0.23 / 1.30 \approx 0.82$. English therefore exposes roughly four-fifths of its grammatical information on the surface form, a high recoverability limited chiefly by irregular forms and by the handful of genuinely ambiguous surface patterns.

The framework's predicted ordering variable for English is $\rho \cdot H \approx 1.07$ bits, second-smallest of the four, predicting a small rebate. \S\ref{sec:results} reports something stronger than a small rebate, namely none: at compute-optimal scale English's grammar-aware model finishes $\Delta\mathrm{bpc} = -0.036$, \emph{worse} than its baseline. English's grammar is legible enough from the surface that a model with sufficient data learns it unaided, and the explicit feature stream only adds cost.

\subsection{Turkish: agglutinative}
\label{sec:cases-tr}

Turkish is the framework's cleanest case. Its word forms are built by stacking suffixes onto a root in a strict slot order: root $\to$ derivational markers $\to$ number $\to$ possessive $\to$ case for nouns, and root $\to$ derivational markers $\to$ voice $\to$ tense/aspect/mood $\to$ person for verbs. Vowel harmony adjusts the surface shape of each suffix to match the vocalic profile of the root, but the harmony rules are regular and invertible; feature tags are harmony-normalised, so \emph{-den} and \emph{-dan} both carry the single tag \texttt{case=ABL}. Morpheme boundaries are unambiguous, and the slot-order constraint rules out combinatorial ambiguity. The engine now performs a holistic verbal and nominal analysis, with deterministic tie-breaks, so that competing analyses of an ambiguous suffix string resolve to a single canonical parse. Two refinements matter for the corpus statistics. A proper-noun fallback intercepts capitalised sentence-initial tokens (\emph{Ahmet}, \emph{\.{I}stanbul}, \emph{T\"urkiye}) and labels them \texttt{pos=PROPN} with an empty bundle, rather than forcing them through the verbal analyser and producing spurious tense and person tags. And denominal-verb derivations in \emph{-le} and \emph{-la\c{s}} (\emph{ba\c{s}-la}, ``to begin''; \emph{g\"uzel-le\c{s}}, ``to become beautiful'') are recognised as a derivational layer before inflectional peeling, so that the verbal bundle attaches to a derived verbal stem rather than to the underlying nominal root.

A representative decomposition. The word \emph{evlerinizden} (``from your (pl.) houses'') resolves to root \emph{ev} (``house'') and feature bundle \texttt{pos=NOUN | num=PL | poss=2PL | case=ABL}, a four-slot stack (plural marker \emph{-ler}, second-person-plural possessive \emph{-iniz}, ablative case \emph{-den}). The word \emph{yaz{\i}ld{\i}} (``it was written'') resolves to root \emph{yaz} (``write'') and bundle \texttt{pos=VERB | voice=PASS | tense=PAST\_DEF}. The surface string is the concatenation of the root and the ordered suffix sequence; for well-formed text, the engine's parse is the unique analysis consistent with the slot-order grammar.

The feature-bundle inventory is large, reflecting the combinatorial nature of Turkish agglutination. The 20{,}000-sentence Turkish sample yields $|B(L)| = 511$ distinct bundles, with $H_{\max} = \log_2 511 \approx 9.00$ bits and corpus-weighted $H(L) = 4.01$ bits. Each slot admits several values, and their Cartesian product yields a large effective space even at moderate per-slot cardinality; a proper-noun fallback keeps the count from inflating further by collapsing the long tail of \emph{Ahmet}, \emph{\.{I}stanbul}, \emph{T\"urkiye}-type capitalised forms into a single \texttt{PROPN} signature rather than forcing each through the verbal analyser. The gap between $H_{\max}$ and $H(L)$ reflects the corpus distribution's concentration on a smaller subset of frequently-used combinations.

Structural recoverability is Turkish's distinguishing property. The conditional entropy $H(b \mid \text{signature})$ on the 20{,}000-sentence sample is $1.03$ bits, yielding $\rho(\text{Turkish}) = 1 - 1.03/4.01 \approx 0.74$. Turkish does not reach $\rho = 1$: vowel-harmony variant classes, a small number of irregular verbal stems, and some derivational homographies leave residual surface ambiguity. But $\rho$ for Turkish sits high, close to English's $0.82$ and well above Arabic's $0.53$, confirming that Turkish wears its grammar on the surface.

The framework's predicted ordering variable for Turkish is $\rho \cdot H \approx 2.98$ bits, second only to Arabic. Turkish was, on the strength of its typology, the language we most expected to gain. A cleanly agglutinative system stacks its grammar in transparent, surface-recoverable suffixes, exactly the configuration the framework rewards, and the pilot bore this out emphatically: Turkish posted by far the largest morphology rebate of the three pilot languages and held it across every rung of the small-scale ladder, even where English's and Arabic's rebates had already inverted (\S\ref{sec:results-pilot}). On typological grounds Turkish was the natural poster case, the language whose structure should make grammar-aware tokenisation pay most.

That expectation does not survive scale. At the compute-optimal, parameter-matched scale of \S\ref{sec:results}, Turkish's rebate is gone: $\Delta\mathrm{bpc} = -0.017$, the grammar-aware model finishing \emph{worse} than its baseline despite carrying more parameters and a larger vocabulary. The very transparency that made Turkish the framework's poster case turns out to be its undoing here: because Turkish wears its grammar openly on the surface, a model given enough data simply learns that morphology for itself, and handing it the same suffix analysis explicitly buys nothing while costing a little. Turkish's fall, from the pilot's largest rebate to none at compute-optimal scale, is the clearest single sign that the Agglutinative Compounding Effect was an artefact of under-training rather than a property of agglutinative structure. The prediction the framework got right was a ranking the pilot confirmed; the prediction it got wrong was that the ranking would mean anything once the baselines were properly trained.

\subsection{Arabic: templatic}
\label{sec:cases-ar}

Arabic's morphological system is organised around two orthogonal axes: a consonantal root carrying semantic content (predominantly three consonants, with a smaller class of quadriliteral roots and a handful of longer ones), and a pattern (\ar{وَزْن}, \emph{wazn}) carrying grammatical content. The root combines, in strict consonantal order, with any of a fixed set of patterns to produce a family of related words, each interleaving the root consonants with pattern-specific vowels and affixes. The root \ar{ك-ت-ب} (\emph{k-t-b}, the writing-related meaning) combines with pattern~I to give \ar{كَتَبَ} (\emph{kataba}, ``he wrote''); with pattern~III to give \ar{كَاتَبَ} (\emph{k\=ataba}, ``he corresponded with''); and with the active-participle pattern of form~I to give \ar{كَاتِب} (\emph{k\=atib}, ``writer''). The same consonantal skeleton runs through all of these; the grammatical identity sits in the pattern.

The engine represents Arabic word forms through a four-part decomposition: the surface string, a feature bundle, the extracted triconsonantal root, and a \emph{wazn-class} label that names the semantic role of the pattern. The wazn-class taxonomy partitions the templatic system into 185 semantic classes, capturing the regularity that each Arabic pattern is associated with a recognisable semantic effect on its root (causative, intensive, reciprocal, seeking, and so on). A representative decomposition. The word \ar{اسْتَخْدَمَ} (\emph{istakhdama}, ``he used'') resolves to root \ar{خ-د-م} (\emph{kh-d-m}, the serving-related meaning), feature bundle \texttt{pos=VERB | form=X}, and wazn-class \texttt{FORM\_X\_SEEKING} (the SEEKING\_ESTIMATIVE / REQUEST role, in which Form~X derives a verb of seeking, requesting, or considering-to-be from its base; \emph{istakhdama} therefore reads literally as ``he sought service from'' before lexicalising to ``he used''). The word \ar{الْمِصْرِيَّة} (\emph{al-mi\d{s}riyya}, ``the Egyptian (feminine)'') resolves to root \ar{مِصْر} (\emph{m-\d{s}-r}, Egypt), bundle \texttt{pos=NOM | def=DEF | role=DERIVED}, and the corresponding nominal wazn-class. Arabic prefixes and suffixes (the definite article \ar{الْـ} (\emph{al-}), person markers on verbs, gender and number agreement markers on nouns) appear alongside the pattern as tag values rather than as independent morphemes. The engine propagates the \texttt{semantic\_role} field into the bundle through a lookup step (Step~B$'$, the wazn-class resolver) that runs after templatic identification. It is this propagation layer that turns the templatic-lexical-surplus prediction from a theoretical claim into a testable one.

The feature-bundle inventory is the largest of the four. On the 20{,}000-sentence sample Arabic yields $|B(L)| = 1{,}925$ observed distinct bundles, with $H_{\max} = \log_2 1925 \approx 10.91$ bits and corpus-weighted $H(L) = 6.51$ bits. The high $H(L)$ reflects Arabic's rich grammatical encoding: gender, number, definiteness, case, person, aspect, voice, templatic form, and wazn-class semantic role all live on a single word form, and their realised distribution in natural text is less concentrated than the other languages'.

Structural recoverability is where Arabic's typological character asserts itself most clearly. The observed conditional entropy $H(b \mid \text{signature})$ is $3.04$ bits, far higher than English's $0.23$ or Turkish's $1.03$. This follows from the inherent ambiguity of the templatic system: the engine's template labels collapse many distinct feature bundles into the same surface pattern. For instance, \texttt{VERB\_TRILATERAL\_BARE} is carried by forms of many tense, voice, and agreement combinations, which the template alone cannot distinguish without lexical or contextual resolution. $\rho(\text{Arabic}) = 1 - 3.04/6.51 \approx 0.53$. Only about half of Arabic's grammatical information is recoverable from the surface-pattern signature; the other half requires lexicon access or context. Arabic has, by a wide margin, the lowest $\rho$ of the four languages, the formal statement of the fact that it hides more of its grammar off the surface than any of the others.

The framework's predicted ordering variable for Arabic is $\rho \cdot H \approx 3.47$ bits, the highest of the four. Arabic is the one language for which grammar-aware tokenisation pays at compute-optimal scale: \S\ref{sec:results} records $\Delta\mathrm{bpc} = +0.022$, the only positive rebate of the four. The reason is the low $\rho$ rather than the high $\rho \cdot H$. The half of Arabic's grammar that sits off the surface is precisely what a surface-reading model cannot learn from text on its own, so the explicit root and \emph{wazn} streams supply information that is otherwise out of reach. This is the Templatic Lexical Surplus; whether it is the templatic structure or merely the parameters those streams add that earns the rebate is settled by the ablation of \S\ref{sec:results-ablation}.

